
A practitioner applying CCNet's standard perplexity filter to RedPajama-V2 retains 86.4\% of Spanish documents but only 16.3\% of English documents at a moderate threshold --- not from any quality difference between the two language groups, but from a calibration artifact introduced by applying a global threshold to perplexity scores that were produced by per-language models trained on incommensurable reference corpora. This paper characterizes that artifact precisely: a large, structural, language-family retention divide (Cramér's V = 0.40--0.57) that is consistent across the full practical range of filtering aggressiveness, maps exactly onto the Romance--Germanic divide, and is 25--40\% larger than prior estimates suggested.

Three contributions are reported: (1) the first calibrated Cramér's V baseline for global ccnet\_perplexity thresholding on RedPajama-V2, (2) quantification of the language-family structure of the disparity consistent with the CCNet KenLM mechanism, and (3) recalibration of prior effect size estimates that affects how researchers should design correction studies.

The most urgent next step is testing whether per-language k-th percentile calibration reduces $\Delta V \geq$ 0.10 for $\geq 3$ of 5 k values --- the direct correction hypothesis. If confirmed, per-language calibration should become the default rather than the exception in multilingual dataset curation. The challenge of calibrating multilingual quality filters equitably has been recognized qualitatively for years; this work provides the quantitative foundation needed to evaluate progress against a precise, reproducible standard.

